\chapter{Conclusão e próximos passos}

Em duas horas percorremos o caminho completo de um banco relacional:
\textbf{entender o modelo}, \textbf{subir o servidor}, \textbf{conectar por um
cliente}, \textbf{definir a estrutura} (DDL), \textbf{manipular os dados}
(CRUD) e \textbf{relacionar tabelas} (chave estrangeira e \comando{JOIN}). Isso
é suficiente para modelar e operar o banco de um primeiro projeto da liga.

\section{O que estudar depois}

Nesta ordem:

\begin{enumerate}
    \item \textbf{Normalização (1FN, 2FN, 3FN)} --- como decidir quantas
          tabelas o seu projeto precisa \cite{elmasri2016}.
    \item \textbf{Modelagem entidade-relacionamento} --- desenhar o modelo
          antes de escrever SQL.
    \item \textbf{Relacionamento N:N} com tabela intermediária, e
          \comando{JOIN} de três ou mais tabelas.
    \item \textbf{Agregação avançada}: \comando{GROUP BY}, \comando{HAVING} e
          subconsultas.
    \item \textbf{Índices} --- por que uma consulta fica lenta e como
          acelerá-la (\comando{EXPLAIN}).
    \item \textbf{Transações} --- \comando{START TRANSACTION},
          \comando{COMMIT}, \comando{ROLLBACK} e o conceito ACID.
    \item \textbf{Integração com aplicação} --- conectar via linguagem
          (Python, PHP, Node) e conhecer os ORMs.
    \item \textbf{Backup e restauração} --- \comando{mysqldump}.
\end{enumerate}

\section{Onde praticar sem instalar nada}

\begin{itemize}
    \item \href{https://sqliteonline.com/}{SQLite Online} --- vários bancos direto no navegador
    \item \href{https://www.db-fiddle.com/}{DB Fiddle} --- MySQL e PostgreSQL no navegador
    \item \href{https://onecompiler.com/mysql/}{OneCompiler MySQL}
    \item \href{https://sqlzoo.net/}{SQLZoo} e \href{https://sqlbolt.com/}{SQLBolt} --- exercícios interativos guiados
    \item \href{https://www.hackerrank.com/domains/sql}{HackerRank SQL} --- desafios com correção automática
\end{itemize}

A documentação oficial do MySQL 8.4 \cite{mysql84} é a referência definitiva
para qualquer dúvida de sintaxe.

\section{Materiais desta capacitação}

\begin{itemize}
    \item \comando{docs/} --- este guia e a apresentação
    \item \comando{materials/liga.sql} --- o script completo da aula, pronto
          para colar no phpMyAdmin
    \item \comando{infrastructure/} --- \comando{compose.yml},
          \comando{.env.example} e os scripts de inicialização do banco
\end{itemize}
